\documentclass[10pt,a4paper]{article}

\usepackage{fontspec}
\usepackage{polyglossia}
\setdefaultlanguage{polish}
\setmainfont{DejaVu Sans}
\setsansfont{DejaVu Sans}
\setmonofont{DejaVu Sans Mono}

\usepackage[a4paper,margin=15mm,headheight=15pt]{geometry}
\usepackage{tikz}
\usetikzlibrary{arrows.meta,positioning,fit,calc,backgrounds}
\usepackage{booktabs}
\usepackage{tabularx}
\usepackage{enumitem}
\usepackage{titlesec}
\usepackage{xcolor}
\usepackage{fancyhdr}
\usepackage{pdflscape}
\usepackage[hidelinks]{hyperref}

\definecolor{navy}{HTML}{17324D}
\definecolor{blue}{HTML}{2B6F9F}
\definecolor{cyan}{HTML}{DCEFF7}
\definecolor{green}{HTML}{DDEFE5}
\definecolor{orange}{HTML}{F8E8CF}
\definecolor{redsoft}{HTML}{F5DDDD}
\definecolor{grayline}{HTML}{AEB8C2}
\definecolor{textgray}{HTML}{38444F}

\pagestyle{fancy}
\fancyhf{}
\lhead{Bezpieczny chatbot RAG - dokumentacja UML}
\rhead{Uniwersytet DSW Ideis}
\cfoot{\thepage}
\renewcommand{\headrulewidth}{0.4pt}

\setlength{\parindent}{0pt}
\setlength{\parskip}{3pt}
\setlist[itemize]{leftmargin=5mm,itemsep=1pt,topsep=2pt}
\titlespacing*{\section}{0pt}{1.1ex}{0.5ex}

\tikzset{
  flow/.style={draw=navy, rounded corners=2mm, fill=cyan, minimum height=10mm,
    align=center, font=\small\sffamily, inner xsep=3mm},
  policy/.style={draw=navy, rounded corners=2mm, fill=orange, minimum height=10mm,
    align=center, font=\small\sffamily, inner xsep=3mm},
  arrow/.style={-{Latex[length=2.3mm]}, semithick, draw=navy},
  umlclass/.style={draw=navy, rounded corners=1.5mm, fill=white, align=left,
    font=\sffamily, inner sep=0pt, text width=42mm},
  relation/.style={-{Latex[length=2.2mm]}, semithick, draw=blue},
  compose/.style={-{Latex[length=2.2mm]}, semithick, draw=navy},
  seqhead/.style={draw=navy, rounded corners=1mm, fill=cyan, minimum width=24mm,
    minimum height=9mm, align=center, font=\scriptsize\bfseries\sffamily},
  message/.style={-{Latex[length=1.8mm]}, thin, draw=navy},
  return/.style={-{Latex[length=1.8mm]}, thin, dashed, draw=blue}
}

\newcommand{\classbox}[3]{%
  \begin{minipage}{42mm}
  \centering\colorbox{navy}{\parbox{38mm}{\centering\color{white}\bfseries #1}}\\[-1pt]
  \raggedright\scriptsize\textcolor{textgray}{#2}\\[-1pt]
  \rule{42mm}{0.35pt}\\[-4pt]
  \raggedright\scriptsize #3
  \end{minipage}%
}

\begin{document}

\begin{center}
  {\color{navy}\sffamily\bfseries\LARGE Bezpieczny chatbot RAG w Telegramie}\par
  \vspace{2mm}
  {\large Dokumentacja UML i krótkie omówienie projektu}\par
  \vspace{1mm}
  {\small Temat: asystent regulaminu studiów - numer indeksu 13897 - aktualizacja 15.07.2026}
\end{center}

\vspace{2mm}
\section*{Cel i zakres}

System jest lokalnym asystentem regulaminu studiów działającym w Telegramie. Korzysta z
dokumentu PDF oraz z kuratorowanego wyciągu z oficjalnej strony Uniwersytetu DSW Ideis
Kraków. Wiadomość użytkownika jest normalizowana, filtrowana pod kątem danych wrażliwych,
klasyfikowana przez Bielik Guard i analizowana metodami NER. Następnie hybrydowy retriever
porównuje zapytanie z embeddingami zapisanymi w SQLite. Baza zawiera 59 fragmentów.

Telegram przechowuje maksymalnie trzy ostatnie, już zamaskowane wiadomości. Historia jest
używana wyłącznie dla krótkich dopowiedzeń, takich jak ``ile mam czasu?'' lub ``jak się
odwołać?''. Jednoznaczne fakty mogą zostać zbudowane przez politykę deterministyczną bez
wywołania modelu generatywnego. W pozostałych przypadkach odpowiedź przygotowuje lokalny LLM
uruchamiany przez Ollamę, a wynik przechodzi kontrolę ugruntowania i filtry wyjściowe.

\begin{center}
\begin{tikzpicture}[node distance=3.2mm,scale=0.90,transform shape]
  \node[flow] (tg) {Telegram +\\zamaskowana historia};
  \node[flow, right=of tg] (privacy) {Privacy Filter\\maskowanie PII};
  \node[flow, right=of privacy] (guard) {Bielik Guard\\wejście};
  \node[flow, right=of guard] (nlp) {NER i bramka\\domenowa};
  \node[flow, below=7mm of nlp] (rag) {Retriever hybrydowy\\SQLite + embeddingi};
  \node[policy, left=of rag] (facts) {FactPolicy\\odpowiedź dowodowa};
  \node[flow, left=of facts] (llm) {Lokalny LLM\\Ollama - gdy potrzebny};
  \node[flow, left=of llm] (output) {Grounding + filtry\\wyjściowe};
  \draw[arrow] (tg) -- (privacy);
  \draw[arrow] (privacy) -- (guard);
  \draw[arrow] (guard) -- (nlp);
  \draw[arrow] (nlp) -- (rag);
  \draw[arrow] (rag) -- (facts);
  \draw[arrow] (facts) -- node[above,font=\scriptsize]{fallback} (llm);
  \draw[arrow] (facts.south) -- ++(0,-4mm) -| (output.south);
  \draw[arrow] (llm) -- (output);
\end{tikzpicture}
\end{center}

\section*{Komponenty i źródła wiedzy}

\begin{tabularx}{\textwidth}{@{}>{\bfseries\color{navy}\raggedright\arraybackslash}p{43mm}X@{}}
\toprule
Komponent & Odpowiedzialność \\
\midrule
SecureRAGPipeline & Koordynuje filtry, kontekst rozmowy, retrieval, generację, grounding i audyt. \\
PrivacyFilter + SafetyClassifier & Maskują PII oraz klasyfikują wejście i wyjście przez Bielik Guard. \\
VectorStore + Embedder & Przechowują 59 fragmentów i realizują ranking embeddingowy oraz BM25. \\
FactPolicy & Buduje krótkie odpowiedzi dla faktów jednoznacznie związanych ze źródłem. \\
OllamaLLM & Generuje odpowiedź tylko wtedy, gdy nie wystarcza bezpieczna reguła dowodowa. \\
Ingest PDF/TEXT & Indeksuje PDF oraz pliki MD/TXT będące kontrolowanymi wyciągami ze stron WWW. \\
AuditLogger & Zapisuje statusy, encje, czasy i źródła bez surowych danych wrażliwych. \\
\bottomrule
\end{tabularx}

\vspace{2mm}
\textbf{Źródła:} regulamin studiów IDEIS 2025/2026 oraz wyciąg z oficjalnych stron
\url{https://www.ideis.pl/krakow/} i \url{https://www.ideis.pl/krakow/kontakt}.

\textbf{Weryfikacja:} 31/31 testów jednostkowych, 40/40 przypadków granicznych i 11/11
testów faktograficznych z lokalnym RAG. Testy obejmują m.in. PII, prompt injection,
literówki, pytania spoza domeny, ugruntowanie liczb, kontekst dopowiedzeń i dane Krakowa.

\begin{landscape}
{\Large\bfseries UML - diagram klas i zależności\par}
\vspace{4mm}

\begin{center}
\begin{tikzpicture}[x=1cm,y=1cm]
  \node[umlclass] (telegram) at (0,5.8) {\classbox{TelegramApp}
    {pipeline, chat\_data.safe\_history}
    {+ message(update, context)\\+ start(update, context)}};

  \node[umlclass, text width=50mm] (pipeline) at (6.5,5.8) {\begin{minipage}{50mm}
    \centering\colorbox{navy}{\parbox{46mm}{\centering\color{white}\bfseries SecureRAGPipeline}}\\[-1pt]
    \raggedright\scriptsize\textcolor{textgray}{settings, store, privacy, safety, nlp, llm, audit}\\[-1pt]
    \rule{50mm}{0.35pt}\\[-4pt]
    \raggedright\scriptsize + handle(text, conversation\_context): BotResponse\\
    + is\_followup\_query(text, context): bool
    \end{minipage}};

  \node[umlclass] (response) at (13,5.8) {\classbox{BotResponse}
    {text, status, sources, timings, backend}
    {+ debug\_dict(): dict}};

  \node[umlclass] (privacy) at (0,1.2) {\classbox{PrivacyFilter}
    {model, regex, mask policy}
    {+ analyze(text): PrivacyResult}};

  \node[umlclass] (safety) at (4.8,1.2) {\classbox{SafetyClassifier}
    {Bielik Guard, threshold, fallback}
    {+ analyze(text): SafetyResult}};

  \node[umlclass] (nlp) at (9.6,1.2) {\classbox{NLPAnalyzer}
    {encje, kategoria, słowa kluczowe}
    {+ analyze(text): NLPResult}};

  \node[umlclass] (facts) at (14.4,1.2) {\classbox{FactPolicy}
    {reguły dowodowe, cytowania}
    {+ answer(question, chunks)}};

  \node[umlclass] (store) at (19.2,1.2) {\classbox{VectorStore}
    {SQLite, embedder, blokada}
    {+ search(query, top\_k)\\+ add\_many(records)}};

  \node[umlclass] (llm) at (14.6,-3.4) {\classbox{OllamaLLM}
    {url, model}
    {+ generate(question, chunks)}};

  \node[umlclass] (ingest) at (0,-3.4) {\classbox{Ingestor}
    {PDF, MD, TXT, podział na fragmenty}
    {+ ingest\_pdf(path)\\+ ingest\_text(path)\\+ ingest\_path(path)}};

  \node[umlclass] (embedder) at (7.3,-3.4) {\classbox{Embedder}
    {Ollama, normalized-v1, hash fallback}
    {+ embed\_many(texts)\\+ hash\_embed(text)}};

  \node[umlclass] (audit) at (19.4,-3.4) {\classbox{AuditLogger}
    {JSONL, lock, masked input}
    {+ write(response, masked\_input)}};

  \draw[relation] (telegram) -- (pipeline);
  \draw[relation] (pipeline) -- node[above,font=\scriptsize]{zwraca} (response);
  \draw[compose] (pipeline.south west) -- (privacy.north east);
  \draw[compose] (pipeline.south west) -- (safety.north);
  \draw[compose] (pipeline.south) -- (nlp.north);
  \draw[compose] (pipeline.south east) -- (facts.north west);
  \draw[compose] (pipeline.south east) -- (store.north west);
  \draw[compose] (pipeline.south east) -- ++(0,-1.0) -| (llm.north);
  \draw[compose] (pipeline.south east) -- ++(0,-1.0) -| (audit.north);
  \draw[relation] (ingest) -- node[above,font=\scriptsize]{zasila} (store);
  \draw[compose] (store.south) -- (embedder.north east);

  \node[draw=grayline,dashed,rounded corners,fit=(privacy)(safety)(nlp)(facts)(store),
    inner sep=3mm] {};
\end{tikzpicture}
\end{center}

\vfill
\small
\textbf{Interpretacja.} Telegram przekazuje do fasady bieżącą wiadomość i maksymalnie trzy
zamaskowane wpisy historii. Pipeline posiada komponenty bezpieczeństwa, NLP, retrieval,
politykę faktograficzną i LLM. Ingestor obsługuje zarówno PDF, jak i kontrolowane źródła
tekstowe. VectorStore przechowuje oryginalny tekst, metadane źródła i embedding, natomiast
AuditLogger otrzymuje wyłącznie zamaskowane dane.

\newpage
{\Large\bfseries UML - diagram sekwencji obsługi wiadomości\par}
\vspace{4mm}

\begin{center}
\begin{tikzpicture}[x=1cm,y=1cm]
  \def\top{6.4}
  \def\bottom{-6.0}
  \node[seqhead] (user) at (0,\top) {Użytkownik};
  \node[seqhead] (telegram) at (2.75,\top) {Telegram Bot};
  \node[seqhead] (pipe) at (5.5,\top) {Pipeline};
  \node[seqhead] (privacy) at (8.25,\top) {Privacy Filter};
  \node[seqhead] (guard) at (11,\top) {Bielik Guard};
  \node[seqhead] (rag) at (13.75,\top) {VectorStore};
  \node[seqhead] (facts) at (16.5,\top) {FactPolicy};
  \node[seqhead] (ollama) at (19.25,\top) {Ollama LLM};
  \node[seqhead] (audit) at (22,\top) {AuditLogger};

  \foreach \n in {user,telegram,pipe,privacy,guard,rag,facts,ollama,audit}
    \draw[dashed,draw=grayline] (\n.south) -- (\n.south |- 0,\bottom);

  \draw[message] (user.south |- 0,5.5) -- node[above,font=\scriptsize]{wiadomość} (telegram.south |- 0,5.5);
  \node[fill=green,rounded corners,font=\scriptsize,inner sep=1.5mm] at (2.75,4.85)
    {wybór do 3 zamaskowanych wpisów historii};
  \draw[message] (telegram.south |- 0,4.2) -- node[above,font=\scriptsize]{handle(text, context)} (pipe.south |- 0,4.2);

  \draw[message] (pipe.south |- 0,3.5) -- node[above,font=\scriptsize]{analyze(text)} (privacy.south |- 0,3.5);
  \draw[return] (privacy.south |- 0,3.0) -- node[above,font=\scriptsize]{masked text + PII} (pipe.south |- 0,3.0);
  \draw[message] (pipe.south |- 0,2.35) -- node[above,font=\scriptsize]{analyze(masked text)} (guard.south |- 0,2.35);
  \draw[return] (guard.south |- 0,1.85) -- node[above,font=\scriptsize]{category + confidence} (pipe.south |- 0,1.85);

  \node[fill=orange,rounded corners,font=\scriptsize,inner sep=1.5mm] at (5.5,1.15)
    {NER, domena, follow-up i normalizacja literówek};
  \draw[message] (pipe.south |- 0,0.45) -- node[above,font=\scriptsize]{search(expanded query, 8)} (rag.south |- 0,0.45);
  \draw[return] (rag.south |- 0,-0.05) -- node[above,font=\scriptsize]{PDF/TEXT candidates} (pipe.south |- 0,-0.05);
  \draw[message] (pipe.south |- 0,-0.75) -- node[above,font=\scriptsize]{answer(query, candidates)} (facts.south |- 0,-0.75);
  \draw[return] (facts.south |- 0,-1.25) -- node[above,font=\scriptsize]{answer lub brak reguły} (pipe.south |- 0,-1.25);

  \node[fill=green,rounded corners,font=\scriptsize,inner sep=1.5mm] at (11,-2.0)
    {jeżeli odpowiedź dowodowa istnieje - pomiń LLM};
  \draw[message] (pipe.south |- 0,-2.7) -- node[above,font=\scriptsize]{prompt + top-3 chunks} (ollama.south |- 0,-2.7);
  \draw[return] (ollama.south |- 0,-3.2) -- node[above,font=\scriptsize]{candidate answer} (pipe.south |- 0,-3.2);

  \node[fill=redsoft,rounded corners,font=\scriptsize,inner sep=1.5mm] at (8.25,-3.95)
    {grounding, Privacy Filter wyjścia, Bielik Guard wyjścia};
  \draw[message] (pipe.south |- 0,-4.7) -- node[above,font=\scriptsize]{zamaskowane metadane i źródła} (audit.south |- 0,-4.7);
  \draw[return] (pipe.south |- 0,-5.25) -- node[above,font=\scriptsize]{BotResponse + cytowania} (telegram.south |- 0,-5.25);
  \draw[return] (telegram.south |- 0,-5.75) -- node[above,font=\scriptsize]{odpowiedź} (user.south |- 0,-5.75);
\end{tikzpicture}
\end{center}

\vfill
\small
Ścieżki blokujące kończą przetwarzanie przed retrievalem albo generacją. Dotyczy to m.in.
zablokowanego PII, treści niebezpiecznej, prompt injection i pytań spoza domeny. Jeżeli
źródło nie wiąże podmiotu pytania z konkretną liczbą lub procedurą, system zwraca
\texttt{insufficient\_context} zamiast łączyć niezależne fragmenty. Odpowiedź deterministyczna
i odpowiedź LLM podlegają filtrom wyjściowym, a log audytowy nie zawiera surowej wiadomości.
\end{landscape}

\end{document}
